\section{2024-12-17}
\page

\startframedtext[width=\dimexpr\textwidth+\rightmarginwidth+\rightmargindistance\relax,frame=off,offset=none] % This spans textwidth + margin width
\startplacetable[location={here},title={}]
\setupTABLE[c][1][background=color,backgroundcolor=gray]
\setupTABLE[r][1][background=color,backgroundcolor=gray,style=bold]
\bTABLE[offset=1mm,option=stretch]
\startCSV
Zitat,1,2,3,4,5,6
Öffentlichkeit,Aussage richtet sich an breite Öffentlichkeit; medial verbreitet,In Talkshow,Wahlkampfveranstaltungen vor großem Publikum,Fernsehansprachen der Staatsoberhäupter,Wahlprogramm richtet sich an Wähler und Öffentlichkeit,Rede im Bundestag; öffentliches politisches Forum
Massenmedialität,,,"Aussagen medial verbreitet (Presse, Nachrichten)",Internationale Verbreitung durch Medien,,
Besondere mdl. Kommunikationsformen,,Talkshow „Markus Lanz“; medialer Diskurs,,Offizielle Fernsehansprachen,,Parlamentarische Rede im Bundestag
Repräsentativ,,,,,"Position der FDP als Partei, Vertreter der Wähler",Vertreterin der Grünen vertritt Fraktionsposition
Situationsgebundenheit,Politischer Kontext; Stimmungsmache gegen Gruppen,Benzinpreisentwicklung und soziale Ungleichheit,Wahlkampfsituation; Angriff auf Gegner und Einwanderer,Beginn des Ukraine-Krieges; politische Rechtfertigungen,Vor Bundestagswahl 2021; Abgrenzung von extremen Strömungen,Debatte zur Einführung des Lobbyregisters
Prozessualität,,Politische Diskussion zu Energiekosten und Alternativen,,,Verknüpft mit Wahlprozess und politischen Reformen,Teil des Reformprozesses zur Verbesserung der Transparenz
Persuasion,Aggressive Rhetorik zur Diskreditierung bestimmter Gruppen,,Polarisierende und emotionalisierte Sprache,Politische Rechtfertigung (Putin) und Kritik (Biden),Rhetorik zur Positionierung gegen extreme Ideologien,
\stopCSV
\eTABLE
\stopplacetable
\stopframedtext
